\documentclass{beamer}
\usetheme{Madrid}
\usecolortheme{default}

% Packages
\usepackage{graphicx}
\usepackage{amsmath, amssymb}
\usepackage{booktabs}
\usepackage{listings}

% Title Information
\title{ACAMS Exam Study Guide}
\subtitle{Key AML/CFT Concepts from Deepseek Claude Questions}
\author{Compliance Training}
\date{\today}

\begin{document}
	
	% --- SLIDE 1: TITLE ---
	\begin{frame}
		\titlepage
	\end{frame}
	
	% --- SLIDE 2: GLOBAL FRAMEWORK ---
	\begin{frame}{What is AML/CFT?}
		\begin{block}{Anti-Money Laundering (AML)}
			Laws and procedures to prevent criminals from disguising illegally obtained funds as legitimate income.
		\end{block}
		\begin{block}{Counter-Financing of Terrorism (CFT)}
			Efforts to prevent the financial support of terrorist organizations and activities.
		\end{block}
		\textbf{Goal:} Protect the integrity of the global financial system.
	\end{frame}
	
	% --- SLIDE 3: FATF ---
	\begin{frame}{Key Players: FATF}
		\begin{columns}[T]
			\column{0.5\textwidth}
			\begin{alertblock}{Financial Action Task Force}
				\begin{itemize}
					\item Global standard-setter for AML/CFT.
					\item Sets international standards.
					\item Researches ML/TF methods.
					\item Assesses country effectiveness.
				\end{itemize}
			\end{alertblock}
			\column{0.5\textwidth}
			\begin{exampleblock}{FATF Lists}
				\begin{itemize}
					\item \textbf{Grey List:} Jurisdictions under increased monitoring.
					\item \textbf{Black List:} High-risk jurisdictions subject to call for action.
				\end{itemize}
			\end{exampleblock}
		\end{columns}
	\end{frame}
	
	% --- SLIDE 4: FATF IMMEDIATE OUTCOMES ---
	\begin{frame}{FATF Immediate Outcomes}
		\begin{definition}[Immediate Outcome 3]
			\textbf{Financial supervision} is risk-based and effective.
		\end{definition}
		\begin{definition}[Immediate Outcome 4]
			\textbf{Financial institutions} apply preventive measures consistent with their ML/TF risks.
		\end{definition}
		\begin{block}{Exam Focus}
			Supervisors prioritize IO.3 (risk-based supervision) and IO.4 (institutions applying preventive measures) to evaluate bank compliance.
		\end{block}
	\end{frame}
	
	% --- SLIDE 5: DORA ---
	\begin{frame}{Digital Operational Resilience Act (DORA)}
		\begin{block}{DORA's Regulatory Expectation}
			\begin{center}
				\textbf{Centralizing oversight} of Information and Communications Technology (ICT) risk management across all member-state branches.
			\end{center}
		\end{block}
		\begin{alertblock}{Key Action for Cross-Border Banks}
			Ensure centralized governance of ICT risks rather than relying solely on national regulators' cybersecurity requirements.
		\end{alertblock}
	\end{frame}
	
	% --- SLIDE 6: EARLY EU DIRECTIVES ---
	\begin{frame}{Early European AML Directives}
		\begin{block}{Shortcoming of Early Directives}
			\begin{center}
				\textbf{Failure to ensure timely and complete transposition} by all member states.
			\end{center}
		\end{block}
		\begin{itemize}
			\item Lack of uniform implementation across jurisdictions led to regulatory arbitrage.
			\item Prompted later revisions to impose stricter transposition deadlines and oversight.
		\end{itemize}
	\end{frame}
	
	% --- SLIDE 7: INTERNAL AFC REPORTS ---
	\begin{frame}{Internal AFC Reporting to Governance}
		\begin{block}{Appropriate Contents (Select Three)}
			\begin{enumerate}
				\item Updates on regulatory interactions and supervisory reviews.
				\item Assessment outcomes of control effectiveness on monitoring tools.
				\item Investigation statuses of escalated Suspicious Activity Report (SAR) reviews.
			\end{enumerate}
		\end{block}
		\begin{alertblock}{Note on User Selection}
			Cash flow performance figures are generally \textit{not} a core AFC governance report component.
		\end{alertblock}
	\end{frame}
	
	% --- SLIDE 8: AML POLICY ELEMENTS ---
	\begin{frame}{AML/CFT Policy Components}
		\begin{block}{Key Elements (Select Two)}
			\begin{enumerate}
				\item \textbf{Assignment of compliance responsibilities.}
				\item \textbf{Program review triggers} based on regulatory or business events.
			\end{enumerate}
		\end{block}
		\begin{itemize}
			\item Policies must clearly define roles and the circumstances under which the program is re-evaluated.
		\end{itemize}
	\end{frame}
	
	% --- SLIDE 9: RISK ASSESSMENT (EWRA) ---
	\begin{frame}{Enterprise-Wide Risk Assessment (EWRA)}
		\begin{block}{Scenario: High-Risk Jurisdiction Customers via Mobile}
			The institution must \textbf{incorporate digital delivery channels} into the EWRA to reassess risk exposure.
		\end{block}
		\begin{itemize}
			\item Risk assessments must evolve to include new channels (e.g., mobile platforms) that bypass traditional branch monitoring.
		\end{itemize}
	\end{frame}
	
	% --- SLIDE 10: HIGH-RISK CLIENTS (TCSP) ---
	\begin{frame}{Trust and Company Service Providers (TCSP)}
		\begin{block}{Handling High-Risk Requests}
			\begin{center}
				\textbf{Decline immediately} unless full beneficial owner information and risk justifications are disclosed and verified.
			\end{center}
		\end{block}
		\begin{itemize}
			\item Connections to high-risk industries (NRA) and confidentiality requests are significant red flags.
			\item Offloading responsibility to local counsel is not a valid risk mitigation strategy.
		\end{itemize}
	\end{frame}
	
	% --- SLIDE 11: PEPs (SINGAPORE) ---
	\begin{frame}{Politically Exposed Persons (PEPs)}
		\begin{block}{Best Practice (Singapore AML Framework)}
			Refer the case to the \textbf{high-risk customer review committee} for enhanced due diligence (EDD) before onboarding.
		\end{block}
		\begin{itemize}
			\item PEPs from outside the region require additional scrutiny.
			\item Complex ownership structures and high-value international transactions necessitate committee-level approval.
		\end{itemize}
	\end{frame}
	
	% --- SLIDE 12: RISK FACTORS: MILITARY ---
	\begin{frame}{Military Organizations and Financial Crime}
		\begin{block}{Risk Elements (Select Two)}
			\begin{enumerate}
				\item \textbf{Government ownership} of defense contractors.
				\item \textbf{Oversight by PEPs} (military officials often qualify as PEPs).
			\end{enumerate}
		\end{block}
		\begin{itemize}
			\item These factors increase the risk of bribery, corruption, and diversion of state assets.
		\end{itemize}
	\end{frame}
	
	% --- PART 2: DATA, TECHNOLOGY, AND SCREENING ---
	% --- SLIDE 13: DATA COVERAGE ---
	\begin{frame}{Data Coverage and Gap Assessment}
		\begin{block}{Challenges Identified (Select Two)}
			\begin{enumerate}
				\item \textbf{Missing identifiers in customer profiles.}
				\item \textbf{Inconsistent use of country codes} across systems.
			\end{enumerate}
		\end{block}
		\begin{itemize}
			\item Inconsistencies and missing data severely impact AML controls and monitoring effectiveness.
		\end{itemize}
	\end{frame}
	
	% --- SLIDE 14: DATA CLEANSING ---
	\begin{frame}{Data Cleansing Processes}
		\begin{block}{Examples (Select Two)}
			\begin{enumerate}
				\item \textbf{Removal of duplicate entries.}
				\item \textbf{Validation of data accuracy.}
			\end{enumerate}
		\end{block}
		\begin{itemize}
			\item These actions are critical during post-acquisition integration to address AML control weaknesses.
		\end{itemize}
	\end{frame}
	
	% --- SLIDE 15: VALIDATION AND TESTING ---
	\begin{frame}{AFC Data Validation}
		\begin{block}{Objectives (Select Three)}
			\begin{enumerate}
				\item Ensuring data inputs reflect actual customer and transaction behaviors.
				\item Verifying that automated controls are functioning correctly based on input quality.
				\item Testing whether controls generate alerts under appropriate risk conditions.
			\end{enumerate}
		\end{block}
	\end{frame}
	
	% --- SLIDE 16: AI TOOLS ---
	\begin{frame}{Transitioning to AI-Based Compliance}
		\begin{block}{Maintain Effectiveness (Select Three)}
			\begin{enumerate}
				\item Ensuring \textbf{data quality} for training algorithms.
				\item Managing the change in \textbf{employee workflows}.
				\item Verifying \textbf{transparency and explainability} in AI models.
			\end{enumerate}
		\end{block}
		\begin{itemize}
			\item The correct focus is on data, people, and explainability—not just selecting external auditors.
		\end{itemize}
	\end{frame}
	
	% --- SLIDE 17: MACHINE LEARNING ---
	\begin{frame}{ML Model Validation}
		\begin{block}{Ensuring Regulatory Compliance}
			\begin{center}
				\textbf{Document model development}, ensure explainability, and perform validation \textbf{before deployment}.
			\end{center}
		\end{block}
		\begin{itemize}
			\item Regulators require validation and explainability of all models used in monitoring.
		\end{itemize}
	\end{frame}
	
	% --- SLIDE 18: ENTITY RESOLUTION ---
	\begin{frame}{Blockchain Analytics: Entity Resolution}
		\begin{definition}[Entity Resolution]
			\textbf{Links data from multiple sources} to improve match accuracy.
		\end{definition}
		\begin{itemize}
			\item Critical for identifying entities (e.g., wallets, addresses) across different blockchain networks and data sets.
		\end{itemize}
	\end{frame}
	
	% --- SLIDE 19: FUZZY LOGIC ---
	\begin{frame}{Fuzzy Logic in PEP Screening}
		\begin{block}{Why Integrate Fuzzy Logic?}
			\begin{center}
				Because it accommodates \textbf{inconsistencies in data formats and naming conventions}.
			\end{center}
		\end{block}
		\begin{itemize}
			\item PEP lists often contain variations (e.g., "William" vs. "Bill," misspellings).
			\item Fuzzy logic matches approximate strings, increasing detection rates while reducing false positives.
		\end{itemize}
	\end{frame}
	
	% --- SLIDE 20: PASSIVE LIVENESS ---
	\begin{frame}{Passive vs. Active Liveness Detection}
		\begin{block}{Passive Liveness Detection}
			\begin{center}
				Uses \textbf{minimal user interaction} and analyzes subtle facial movement or texture patterns.
			\end{center}
		\end{block}
		\begin{itemize}
			\item \textbf{Active Methods:} Require user actions (e.g., blinking, smiling, answering KYC questions).
			\item Passive is less intrusive and often preferred for a smoother customer experience.
		\end{itemize}
	\end{frame}
	
	% --- PART 3: TRANSACTION MONITORING & ANALYSIS ---
	% --- SLIDE 21: ROUND TRIP ---
	\begin{frame}{Rules-Based Alerts: "Round Trip"}
		\begin{block}{Analyst Workflow}
			\begin{enumerate}
				\item Review alert (nearly identical send and receive remittance within minutes).
				\item Confirm the transaction pattern.
				\item \textbf{Escalate to Level 2} for deeper investigation.
			\end{enumerate}
		\end{block}
		\begin{itemize}
			\item Level 1 analysts perform initial triage; Level 2 conducts full investigations.
			\item "Round tripping" is a classic layering technique.
		\end{itemize}
	\end{frame}
	
	% --- SLIDE 22: HIGH-VOLUME TRANSACTIONS ---
	\begin{frame}{Card-Based Products: Risk Factors}
		\begin{block}{Common Risk Factors (Select Three)}
			\begin{enumerate}
				\item \textbf{Anonymity} through prepaid cards.
				\item \textbf{High volumes of rapid transactions.}
				\item \textbf{Use by customers in high-risk sectors.}
			\end{enumerate}
		\end{block}
	\end{frame}
	
	% --- SLIDE 23: PREPAID CARDS ---
	\begin{frame}{Prepaid Card Vulnerabilities}
		\begin{block}{Key Vulnerability}
			\begin{center}
				\textbf{Can be used anonymously}, especially when reloadable or unregistered.
			\end{center}
		\end{block}
		\begin{itemize}
			\item This anonymity complicates customer due diligence (CDD) and transaction monitoring.
		\end{itemize}
	\end{frame}
	
	% --- SLIDE 24: E-COMMERCE ---
	\begin{frame}{Transaction Monitoring in E-Commerce}
		\begin{block}{The Challenge}
			\begin{center}
				\textbf{High transaction volumes and decentralized payments} make patterns harder to detect.
			\end{center}
		\end{block}
		\begin{itemize}
			\item E-commerce platforms often process millions of small transactions, making it difficult to separate legitimate from illicit activity.
		\end{itemize}
	\end{frame}
	
	% --- SLIDE 25: REAL ESTATE & PEPS ---
	\begin{frame}{High-Value Assets and Real Estate}
		\begin{block}{Red Flag}
			\begin{center}
				Making payments for high-value items in \textbf{cryptocurrency}.
			\end{center}
		\end{block}
		\begin{itemize}
			\item While complex financing can be a red flag, cryptocurrency payments for high-value assets are increasingly scrutinized due to the potential for anonymity.
			\item A foreign PEP purchasing multiple properties via a trust should trigger an AML investigation.
		\end{itemize}
	\end{frame}
	
	% --- SLIDE 26: FREE-TRADE ZONES (FTZs) ---
	\begin{frame}{Free-Trade Zone Vulnerabilities}
		\begin{block}{Typical Characteristics (Select Two)}
			\begin{enumerate}
				\item \textbf{Limited customs authority} over transitory goods.
				\item \textbf{Evasions of customs payments.}
			\end{enumerate}
		\end{block}
		\begin{itemize}
			\item FTZs are vulnerable to mis-invoicing, over/under-invoicing, and illicit trade.
		\end{itemize}
	\end{frame}
	
	% --- SLIDE 27: PRIVATE BANKING ---
	\begin{frame}{Private Banking Vulnerabilities}
		\begin{block}{Core ML Vulnerability}
			\begin{center}
				\textbf{Discretion and high client autonomy} can obscure beneficial ownership and reduce transparency.
			\end{center}
		\end{block}
	\end{frame}
	
	% --- SLIDE 28: INVESTMENT BANKING ---
	\begin{frame}{Investment Banking Red Flags}
		\begin{block}{Red Flag Scenario}
			Client typically deals in low-margin bond offerings but suddenly issues high-return instruments via an unfamiliar subsidiary.
		\end{block}
		\begin{block}{Prudent Response}
			\begin{center}
				\textbf{Escalate internally} and request further documentation before determining next steps.
			\end{center}
		\end{block}
	\end{frame}
	
	% --- PART 4: REGULATIONS & REPORTING ---
	% --- SLIDE 29: BSAA / FIU ---
	\begin{frame}{Bank Secrecy Act (BSA)}
		\begin{block}{Purpose of FinCEN}
			\begin{center}
				Safeguards the financial system by administering and enforcing the Bank Secrecy Act to \textbf{detect and prevent financial crimes}.
			\end{center}
		\end{block}
	\end{frame}
	
	% --- SLIDE 30: CONSUMER PROTECTION ---
	\begin{frame}{Consumer Protection Objectives}
		\begin{block}{Primary Objectives (Select Two)}
			\begin{enumerate}
				\item \textbf{Prevent unauthorized transactions} affecting vulnerable customers.
				\item \textbf{Ensure clear and transparent disclosures} of terms.
			\end{enumerate}
		\end{block}
	\end{frame}
	
	% --- SLIDE 31: REPORTING INCONSISTENCIES ---
	\begin{frame}{AFC Regulatory Reports}
		\begin{block}{Impact of Inconsistent Reporting}
			\begin{center}
				The institution \textbf{reduces the efficiency of regulatory review}.
			\end{center}
		\end{block}
		\begin{itemize}
			\item Mixed free-text and structured formats hinder supervisors' ability to quickly review data.
			\item Standardization is key to effective compliance.
		\end{itemize}
	\end{frame}
	
	% --- SLIDE 32: DATA PRIVACY ---
	\begin{frame}{Data Privacy and Security}
		\begin{block}{Compliance with Data Privacy Laws}
			\begin{enumerate}
				\item Create \textbf{jurisdiction-specific procedures} for capturing research findings.
				\item Develop a standard research template with \textbf{adjustable sections} based on local data regulations.
			\end{enumerate}
		\end{block}
	\end{frame}
	
	% --- PART 5: SECTORS & SPECIALIZED ROLES ---
	% --- SLIDE 33: MSBS ---
	\begin{frame}{Money Services Businesses (MSBs)}
		\begin{block}{Why Classification Matters}
			\begin{center}
				MSB designation triggers \textbf{specific AML compliance and reporting obligations}.
			\end{center}
		\end{block}
	\end{frame}
	
	% --- SLIDE 34: PSPS ---
	\begin{frame}{Payment Service Providers (PSPs)}
		\begin{block}{Cross-Border Strategy}
			\begin{center}
				\textbf{Customize AML policy and onboarding procedures} per jurisdiction and product.
			\end{center}
		\end{block}
	\end{frame}
	
	% --- SLIDE 35: CASINOS ---
	\begin{frame}{Casino AML Controls}
		\begin{block}{Control Measures (Select Two)}
			\begin{enumerate}
				\item \textbf{Enforcing chip redemption policies} tied to minimum gaming thresholds.
				\item \textbf{Tracking player activity} through registered loyalty programs.
			\end{enumerate}
		\end{block}
		\begin{itemize}
			\item These controls help mitigate the risk of using casinos to layer illicit funds.
		\end{itemize}
	\end{frame}
	
	% --- SLIDE 36: CANNABIS ---
	\begin{frame}{Cannabis-Related Businesses}
		\begin{block}{Federally Regulated Institution}
			\begin{center}
				Apply \textbf{enhanced due diligence}, assess provincial/state law compliance, and consult internal legal guidance.
			\end{center}
		\end{block}
	\end{frame}
	
	% --- SLIDE 37: CRYPTO ---
	\begin{frame}{Cryptocurrency and Sanctions}
		\begin{block}{Minimizing Sanctions Risk}
			Implement robust \textbf{blockchain analytics} and transaction monitoring tools to screen wallet addresses against sanctions lists.
		\end{block}
		\begin{itemize}
			\item Entity resolution links on-chain addresses to off-chain risk data.
		\end{itemize}
	\end{frame}
	
	% --- SLIDE 38: PERPETUAL KYC ---
	\begin{frame}{Perpetual KYC}
		\begin{block}{Technology Solution}
			\begin{center}
				Implement \textbf{perpetual KYC} with centralized data and trigger-based updates.
			\end{center}
		\end{block}
		\begin{itemize}
			\item This reduces customer friction (complaints about repeated KYC requests) while maintaining compliance.
		\end{itemize}
	\end{frame}
	
	% --- SLIDE 39: DATA QUALITY ---
	\begin{frame}{Data Governance and Quality}
		\begin{block}{Common Issues}
			\begin{itemize}
				\item \textbf{Missing identifiers} in customer profiles.
				\item \textbf{Inconsistent address formats} and duplicate profiles.
				\item \textbf{Misalignment} between transaction monitoring scenarios and typologies.
			\end{itemize}
		\end{block}
		\begin{alertblock}{Remedy}
			Initiate a \textbf{data cleansing process} targeting missing and duplicated records.
		\end{alertblock}
	\end{frame}
	
	% --- PART 6: COMPLIANCE CULTURE & STAKEHOLDERS ---
	% --- SLIDE 40: STAKEHOLDER MANAGEMENT ---
	\begin{frame}{Stakeholder Management}
		\begin{block}{Who are the Key Stakeholders?}
			\begin{itemize}
				\item \textbf{Internal:} Board of Directors, Compliance Committee, Business Units.
				\item \textbf{External:} Regulators, FIUs, Law Enforcement, Customers.
			\end{itemize}
		\end{block}
	\end{frame}
	
	% --- SLIDE 41: CULTURE AND TRAINING ---
	\begin{frame}{Compliance Culture}
		\begin{block}{Building a Strong Culture}
			\begin{itemize}
				\item \textbf{Tone from the Top:} Leadership must prioritize compliance.
				\item \textbf{Training:} Regular, role-specific training on AML/CFT policies.
				\item \textbf{Awareness:} Employees should understand the "why" behind controls.
			\end{itemize}
		\end{block}
	\end{frame}
	
	% --- SLIDE 42: POST-ACQUISITION INTEGRATION ---
	\begin{frame}{Post-Acquisition Integration}
		\begin{block}{AML Data Strategy}
			\begin{enumerate}
				\item Conduct a \textbf{data coverage and gap assessment}.
				\item Initiate \textbf{data cleansing} processes.
				\item Harmonize \textbf{country code standards} across merged systems.
			\end{enumerate}
		\end{block}
	\end{frame}
	
	% --- PART 7: FUTURE TRENDS & FINAL REVIEW ---
	% --- SLIDE 43: FUTURE TRENDS ---
	\begin{frame}{Future Trends in AFC}
		\begin{itemize}
			\item \textbf{AI and Machine Learning:} Enhanced detection and prediction.
			\item \textbf{Blockchain Analytics:} Deeper on-chain monitoring.
			\item \textbf{Data Sharing:} Public-private partnerships for information exchange.
			\item \textbf{RegTech:} Automation of compliance and reporting.
		\end{itemize}
	\end{frame}
	
	% --- SLIDE 44: GREY LIST REMEDIATION ---
	\begin{frame}{Remediating FATF Grey List Status}
		\begin{block}{Key Steps}
			\begin{enumerate}
				\item \textbf{Address strategic deficiencies} identified by FATF.
				\item Improve supervision of DNFBPs.
				\item Enhance beneficial ownership transparency.
				\item Demonstrate effective implementation of AML/CFT measures.
			\end{enumerate}
		\end{block}
	\end{frame}
	
	% --- SLIDE 45: KEY TAKEAWAYS ---
	\begin{frame}{Key Takeaways}
		\begin{enumerate}
			\item \textbf{Know Your Customer (KYC)} is the foundation of AML.
			\item \textbf{Risk-based approach} is key: assess, mitigate, monitor.
			\item \textbf{Data quality} is non-negotiable for effective controls.
			\item \textbf{FATF standards} are the global benchmark.
			\item \textbf{Technology} is a tool, not a replacement for human judgment.
		\end{enumerate}
	\end{frame}
	
	% --- SLIDE 46: EXAM TIPS ---
	\begin{frame}{ACAMS Exam Tips}
		\begin{itemize}
			\item \textbf{Read carefully:} Look for words like "Select Three" or "Select Two."
			\item \textbf{Think like a compliance officer:} Escalate, document, and validate.
			\item \textbf{Focus on the "why":} Understand the underlying logic of the risk-based approach.
			\item \textbf{Know the frameworks:} FATF, DORA, Wolfsberg, EU Directives.
		\end{itemize}
	\end{frame}
	
	% --- SLIDE 47: SAMPLE QUESTION 1 ---
	\begin{frame}{Sample Question: Risk Factors}
		\begin{block}{Question}
			Which of the following represent common risk factors associated with card-based products? (Select Three.)
		\end{block}
		\begin{block}{Answer}
			\begin{enumerate}
				\item Anonymity through prepaid cards.
				\item High volumes of rapid transactions.
				\item Use by customers in high-risk sectors.
			\end{enumerate}
		\end{block}
	\end{frame}
	
	% --- SLIDE 48: SAMPLE QUESTION 2 ---
	\begin{frame}{Sample Question: PEPs}
		\begin{block}{Question}
			A compliance officer at a financial institution in Singapore is reviewing a PEP case. What action is most consistent with Singapore’s AML framework?
		\end{block}
		\begin{block}{Answer}
			Refer the case to the high-risk customer review committee for enhanced due diligence before onboarding.
		\end{block}
	\end{frame}
	
	% --- SLIDE 49: SAMPLE QUESTION 3 ---
	\begin{frame}{Sample Question: Data Cleansing}
		\begin{block}{Question}
			Which actions are examples of data cleansing processes? (Select Two.)
		\end{block}
		\begin{block}{Answer}
			\begin{enumerate}
				\item Removal of duplicate entries.
				\item Validation of data accuracy.
			\end{enumerate}
		\end{block}
	\end{frame}
	
	% --- SLIDE 50: SAMPLE QUESTION 4 ---
	\begin{frame}{Sample Question: E-Commerce}
		\begin{block}{Question}
			Why is transaction monitoring particularly challenging in an e-commerce risk context?
		\end{block}
		\begin{block}{Answer}
			High transaction volumes and decentralized payments make patterns harder to detect.
		\end{block}
	\end{frame}
	
	% --- SLIDE 51: REFERENCES ---
	\begin{frame}{References}
		\begin{thebibliography}{99}
			\bibitem{1} Financial Action Task Force (FATF). (n.d.). \textit{Recommendations}.
			\bibitem{2} ACAMS. (2026). \textit{CAMS Exam Study Guide}.
			\bibitem{3} Digital Operational Resilience Act (DORA). (2022). Regulation (EU) 2022/2554.
			\bibitem{4} Wolfsberg Group. (n.d.). \textit{AML Principles}.
		\end{thebibliography}
	\end{frame}
	
	% --- SLIDE 52: QUIZ ---
	\begin{frame}{Quick Quiz}
		\begin{enumerate}
			\item What is the primary purpose of fuzzy logic in PEP screening?
			\item Which FATF Immediate Outcomes do supervisors prioritize for evaluating banks?
			\item What is the best action for a TCSP firm when a client requests confidentiality and is linked to high-risk industries?
		\end{enumerate}
	\end{frame}
	
	% --- SLIDE 53: THANK YOU ---
	\begin{frame}
		\begin{center}
			\Huge Thank You
			\vspace{1cm}
			\Large Good luck on your CAMS Exam!
		\end{center}
	\end{frame}
	
\end{document}